\begin{lemma}[Stable change]
Fix a finite loopless multigraph with its exact edge set, a palette of
$k$ colors, and a distinguished edge $e_0$.  Let $(T,\varphi_n)$ be a
positive-rung extended Tashkinov tree.  Let $\mathcal B$ be a hierarchy
base and let
\[
 T_n=T_{n,0}\subset T_{n,1}\subset\cdots\subset
 T_{n,q}\subset T_{n,q+1}=T
\]
be a good hierarchy $H$ under $\varphi_n$.  Put
\[
 A=T_{n,q}^{*}=\noderef{HierarchyStarredStage}(H,q),\qquad
 D_n=\noderef{TashkinovStoredColors}(\mathcal E,q_n),
 \qquad D_{n,q}=\noderef{HierarchyStoredColorsAt}(H,q).
\]
Assume that $T$ has maximum property, that $V(T)$ is not elementary,
that $\varphi_n$ is proper and leaves exactly $e_0$ uncolored, and that
these data form a \noderef{HierarchyCounterexample} $C$.

Assume, uniformly for every positive-rung maximum ETT on the same graph,
palette, and distinguished edge, with arbitrary certified input and
generating colorings and with rung number smaller than that of $T$, the
six lower-rung assertions of the technical theorem: elementarity,
terminal interchangeability, the two PE stability and stored-pair
assertions, the terminal PE component assertion, the defective-endpoint
order assertion, and the mod-coloring assertion.  Assume also that every
same-rung maximum ETT with a good hierarchy having fewer than $q$
dividers is elementary.  Assume that $C$ is lexicographically minimal
among same-rung hierarchy counterexamples with exactly $q$ dividers, in
the sense of \noderef{HierarchyCounterexampleLess} and
\noderef{HierarchyCounterexampleMeasure}, and assume
\[
                 2n+11\leq |\overline\varphi_n(T_n)|.          \tag{L}
\]

Put $\ell=|T|-|A|$, the number of exact tail additions above $A$, and
\[
 T(y_{\ell-1})=T\big|_{\min\{|A|+(\ell-1),|T|\}}.
\]
Assume (6.6), namely that $V(T(y_{\ell-1}))$ is elementary under
$\varphi_n$.  Let $\alpha\ne\beta$ be colors missing on
$T(y_{\ell-1})$, and put
\[
 i_c=\noderef{HierarchyFirstMissingPosition}(\varphi_n,T,c),
 \qquad v_c=T.\mathit{vertex}(i_c).
\]
Let $z$ be an arbitrary node of the two-color incidence graph, let $Q$
be the full $(\alpha,\beta)$-component through $z$, and put
\[
 \sigma_n=\noderef{TwoColorComponentSwap}
       (G,\varphi_n,\alpha,\beta,z).
\]
Suppose at least one of the following three cases occurs:
\begin{enumerate}
\item $q\geq1$, $\alpha\in\overline\varphi_n(T_{n,q})$, $Q$ is a
possibly singleton \noderef{TwoColorPathComponent}, and the full node
support of $Q$ is disjoint from the full component through
$\operatorname{inl}v_\alpha$;
\item $q=0$, either $\alpha\in\overline\varphi_n(T_n)$ or
$\alpha\in\overline\varphi_n(A)$ with $\alpha,\beta\notin D_n$, $Q$ is
a possibly singleton \noderef{TwoColorPathComponent}, and the full node
support of $Q$ is disjoint from the full component through
$\operatorname{inl}v_\alpha$; or
\item every vertex of $A$ precedes $v_\alpha$,
$v_\alpha\prec v_\beta$, $\alpha,\beta\notin D_{n,q}$, and
\[
 \alpha\notin
 \{c:\exists t\in\operatorname{Fin}|T|,
        i_\alpha\leq t<i_\beta,
        \ \varphi_n(T.\mathit{edge}(t))=c\}.                 \tag{U}
\]
In this third case $Q$ is an arbitrary full component; no path or
component-disjointness condition is imposed.
\end{enumerate}

Then the following common conclusion holds.
\begin{enumerate}
\item The coloring $\sigma_n$ is
$(T_{n,q}^{*},D_n,\varphi_n)$-weakly stable in the exact sense of
\noderef{HierarchyWeaklyStable}.
\item The coloring $\sigma_n$ is a
$\varphi_n\bmod T_n$ coloring in the sense of \noderef{ModColoring}.
Either $A=T_n$, the non-strict endpoint which the tablet's positive-rung
ETT type intentionally does not repackage, or $A$ has an actual maximum
ETT presentation under $\sigma_n$.
\item If $q\geq1$, the same ordered base core and closure and the stages
\[
 T_n=T_{n,0}\subset\cdots\subset T_{n,q}=A
\]
form a good hierarchy of an actual maximum ETT on $A$ under $\sigma_n$.
Its divider count is $q-1$; at every retained index the sets $D_{n,j}$
agree exactly, and $\Gamma^j_\eta$ agrees exactly for every
$\eta\in D_{n,j}$.  In addition,
\[
 T_{n,q}\text{ is }
 \left(\bigcup_{\eta\in D_{n,q}}\Gamma^{q-1}_\eta\right)^{-}
 \text{-closed under }\sigma_n.                              \tag{M}
\]
\end{enumerate}
Furthermore, whenever the complete certificate in Case 3 holds, the
full tree $T$ has an actual corresponding maximum ETT presentation under
$\sigma_n$, and the displayed hierarchy remains good with the same
ordered base, stages, dividers, $D$-sets, and all paper-defined
$\Gamma$-sets, in the precise sense of \noderef{SameGoodHierarchy}.
This last assertion is an implication from the Case 3 certificate; it
does not assert that Cases 1 and 2 are exclusive from Case 3.
\end{lemma}
\begin{proof}
Denote by $\mathcal C_i$ the complete certificate displayed in Case
$i$, and denote by $\mathcal K$ the three-part common conclusion.
The assumed case disjunction is
$\mathcal C_1\vee\mathcal C_2\vee\mathcal C_3$.
The common same-rung minimality hypothesis, along with every other common
ambient hypothesis, is passed unchanged in each application below.

Suppose first that $\mathcal C_1$ holds.  It supplies, in addition to all
the common ambient hypotheses, the positive-divider missing-color
condition, the possibly singleton path-component certificate, and
disjointness of its full incidence-component support from the component
through $v_\alpha$.  These are exactly the first-case inputs of
\noderef{StableChangeCasesOneTwo}.  That lemma therefore gives
$\mathcal K$.

Suppose instead that $\mathcal C_2$ holds.  Its zero-divider alternative
is exactly the second alternative in
\noderef{StableChangeCasesOneTwo}, and it supplies the same path and full
component-disjointness certificates required there.  Applying that lemma
with its second case therefore again gives $\mathcal K$.

Finally suppose that $\mathcal C_3$ holds.  Its four conjuncts are,
respectively, the assertion that $A$ precedes $v_\alpha$, the order
$v_\alpha\prec v_\beta$, exclusion of both colors from $D_{n,q}$, and
the exact segment-unused assertion (U).  Together with the shared
ambient hypotheses and the arbitrary component through $z$, these are
exactly the inputs of \noderef{StableChangeCaseThree}.  Its first four
conclusion fields are precisely $\mathcal K$, so $\mathcal K$ follows in
this case as well.  Dispatching the assumed three-way disjunction through
these three applications proves the common conclusion.

It remains to prove the stated implication.  Assume the complete
certificate $\mathcal C_3$, independently of which disjunct was used for
the preceding common conclusion.  Apply
\noderef{StableChangeCaseThree} once more to this certificate.  Its final
conclusion field is an actual corresponding maximum ETT on $T$, together
with a maximum-property witness and a good hierarchy related to the
original one by the full \noderef{SameGoodHierarchy} relation.  This is
exactly the asserted full-tree conclusion.  Since this argument assumes
$\mathcal C_3$ directly, it requires no exclusivity among the three case
predicates and completes the proof.
\end{proof}
